\section{理想主义与商业化：开路的人}

最后回到价值判断：上市公司要赚钱，但张鹏仍把 AGI 和行业贡献放在更高位置。这里是整期访谈从公司史回到长期理想的地方。

访谈结尾，张鹏被问到做一家实现 AGI 的公司，还是做一家利润很高的公司，二选一选哪个。他选择 AGI，但也强调二者并不对立：如果实现 AGI，也会是一家商业上伟大的公司。更重要的是，他不满意只赚钱而没有技术产出、没有行业贡献的智谱。

\begin{figure}[H]
\centering
\includegraphics[width=0.96\textwidth]{idealism-commercialization.png}
\caption{理想主义与商业化：技术理想、工程文化、客户价值、治理和先行者身份。}
\end{figure}

\begin{knowledgebox}{读图：智谱想要的是平衡路线}
图中从技术理想到客户价值再到上市治理，最后回到“先行者”。张鹏强调的不是空喊 AGI，也不是只做赚钱机器，而是把“让机器像人类一样思考”变成能服务客户和社会的工程系统。
\end{knowledgebox}

\subsection{让机器像人类一样思考}

理想主义不能只停在“AGI”这个词上。本节把智谱的 slogan 拆开，看它如何同时指向技术目标和社会价值。

张鹏用公司 slogan 解释智谱对 AGI 的理解：让机器像人类一样思考，并反过来赋能人类社会。这句话可以拆成两部分：第一是技术目标，机器要具备更强理解、推理和行动能力；第二是价值目标，技术要产生实际价值，而不是停在实验室里。

\begin{importantbox}{AGI 的可操作化}
如果 AGI 定义彼此不同，判断一家公司是否追求 AGI，不能只听口号，而要看它把长期目标拆成哪些阶段任务、技术路线、产品路径和客户价值。
\end{importantbox}

\subsection{先行者而非只做明星项目}

前面讲的是“做什么”，这里讲“希望被怎样记住”。先行者这个词，比明星公司更朴素，也更符合智谱的自我叙事。

节目开头提到，有人评价智谱像“水泥”：能干漂亮的活，但情绪价值不强。张鹏最后说，希望智谱被历史记为 AGI 的先行者、开路的人。这个自我叙事和“水泥”并不冲突：开路的人未必最会制造话题，但要长期铺路、踩坑、转化研究、服务客户、承受组织治理。

\subsection{本章小结}

智谱的自我定位，是在理想主义和商业化之间走平衡路线。上市不是终点，而是让一家追求 AGI 的公司进入更高治理要求下继续前进。

